We formulate the Rivero--de Vries electroweak observation as a two-branch secular problem built from a spin-Casimir variable \(J=s(s+1)\).  The positive root of \(x^2+Jx-J=0\), evaluated at \(J=3/4\) and \(J=2\), gives \(1-x_+(3/4)/x_+(2)=0.22310132\ldots\), close to the weak angle defined by the W/Z pole-mass ratio.  The paper treats this agreement as a pole-spectrum statement and separates it from running weak mixing angles.  We present the algebraic construction, the electroweak projective interpretation in which the Higgs vacuum scale is radial, and the analytical tasks required for a dynamical derivation.  String, brane, Kaluza--Klein, and \(G_2\)-holonomy ingredients are organized as possible mechanisms for the secular equation rather than as independent assumptions.  The central open calculation is the derivation of the ordered assignment \((J_W,J_Z)=(3/4,2)\) from the gauge-Higgs sector or from a compactification boundary problem.
